\documentclass{article}

% HKU COMP3314 course report. Not a NeurIPS 2026 conference submission.
\usepackage[preprint]{neurips_2026}

\usepackage[utf8]{inputenc}
\usepackage[T1]{fontenc}
\usepackage{hyperref}
\usepackage{url}
\usepackage{booktabs}
\usepackage{amsfonts}
\usepackage{amsmath}
\usepackage{amssymb}
\usepackage{microtype}
\usepackage{xcolor}

\title{Spectral Shrunk Clustering: A Critical Study and Reproduction}

\author{%
  Group 23 \\
  Department of Computer Science \\
  The University of Hong Kong \\
  COMP3314 Machine Learning
}

\begin{document}

\maketitle

\begin{abstract}
Spectral clustering has proved effective for finding clusters with complex geometric structure, but its performance depends heavily on how similarities between data points are represented. In \emph{A Convex Formulation for Spectral Shrunk Clustering}, Chang et al.\ (2015) address this issue by learning a refined representation from the spectral embedding rather than relying solely on the neighborhood relationships established in the original feature space. Their method, Spectral Shrunk Clustering (SSC), introduces a convex optimization problem that encourages nearby embedded points to move together while preserving their original positions.

This report examines the motivation, mathematical formulation, optimization procedure, and experimental evidence presented in the paper. We also independently implement SSC and evaluate it on four publicly recoverable benchmark datasets. Although our numerical checks support the objective-consistent optimization formulation, the reported improvements over standard spectral clustering are not consistently reproduced. Further analysis suggests that the discrepancy cannot be attributed to the shrinkage solver alone: the initial spectral embedding, graph construction, and final K-means step also play important roles. These findings highlight the distinction between solving a convex representation-learning problem and obtaining reliable clustering assignments.
\end{abstract}

\section{Introduction}

Clustering aims to organize unlabeled data into groups whose members share similar characteristics. For high-dimensional data, however, the choice of representation can be as important as the clustering algorithm itself. Distances measured directly in the original feature space may not reflect the structure that separates meaningful groups, particularly when the data lie near a lower-dimensional manifold.

Spectral clustering addresses this difficulty by constructing a similarity graph and using the eigenvectors of its graph Laplacian to represent the data in a lower-dimensional space [Shi and Malik, 2000; von Luxburg, 2007]. K-means is then applied to the resulting embedding. Unlike approaches based only on Euclidean distance in the original space, spectral methods can capture more complicated neighborhood relationships.

Nevertheless, the similarity graph is usually constructed before the low-dimensional representation is obtained. This creates a potential mismatch: relationships that appear meaningful in the original feature space may not be equally informative after spectral embedding. Conversely, points that should belong to the same cluster may remain separated in the embedding.

Chang et al.\ (2015) propose Spectral Shrunk Clustering (SSC) to address this mismatch. Rather than directly clustering the initial spectral representation, SSC constructs another neighborhood graph in the embedding space and learns a new representation in which neighboring points are encouraged to move closer together.

The paper makes three principal contributions. First, the refined representation is no longer subject to the orthogonality constraint associated with the spectral eigenvectors, allowing greater flexibility in how the embedded points are positioned. Second, the additional neighborhood graph provides a way to exploit local geometry after dimensionality reduction. Third, the authors formulate the refinement step as a convex optimization problem and propose an iterative solution with a convergence argument.

We examine these claims through both mathematical analysis and experimental reproduction, with particular attention to how the learned representation affects the final clustering results.

\section{Spectral Shrunk Clustering}

\subsection{Spectral embedding and neighborhood graphs}

Let \(X=\{x_1,\ldots,x_n\}\) denote a dataset containing \(n\) observations. Spectral clustering first constructs an affinity matrix \(A\), where \(A_{ij}\) measures the similarity between observations \(x_i\) and \(x_j\).

The paper uses a Gaussian affinity for neighboring points and defines the symmetric normalized graph Laplacian as
\[
L_n=I-D^{-1/2}AD^{-1/2},
\]
where \(D\) is the diagonal degree matrix, with \(D_{ii}=\sum_j A_{ij}\).

The initial spectral representation \(F\in\mathbb{R}^{n\times c}\) is obtained from the \(c\) eigenvectors associated with the smallest eigenvalues of \(L_n\), where \(c\) is the number of clusters.

In conventional spectral clustering, the rows of \(F\) are subsequently partitioned by K-means. SSC instead treats \(F\) as an intermediate representation. A second affinity matrix \(W\) is constructed using the neighborhood relationships among its rows, allowing the algorithm to model local geometry directly in the embedded space.

This distinction is central to the method: \(A\) describes relationships in the original feature space, while \(W\) describes relationships after spectral embedding.

\subsection{The shrinkage objective}

SSC learns a refined representation \(G\in\mathbb{R}^{n\times c}\) by solving
\[
\min_G\;
\|G-F\|_{2,1}
+
\gamma\sum_{i,j}W_{ij}\|g_i-g_j\|_2.
\]

The objective balances two requirements.

The first term, \(\|G-F\|_{2,1}\), discourages the refined representation from moving too far from the original spectral embedding. Because it sums the Euclidean norms of row-wise residuals, it is less sensitive to large individual deviations than a squared Frobenius-norm penalty [Nie et al., 2010].

The second term encourages pairs of points connected by the graph \(W\) to become closer in the refined representation. When \(W_{ij}\) is large, separating \(g_i\) and \(g_j\) incurs a larger penalty.

The parameter \(\gamma\) determines the relative strength of these two effects. A small value keeps \(G\) close to \(F\), whereas a larger value encourages stronger shrinkage along graph edges. Excessive shrinkage may remove distinctions that are useful for clustering, so the choice of \(\gamma\) is consequential.

For fixed \(F\) and nonnegative \(W\), both terms are convex in \(G\). Unlike the spectral embedding problem, the new representation is not required to satisfy an orthogonality constraint. The resulting optimization therefore has a useful convex structure, although its solution is not necessarily unique.

\subsection{Optimization and final cluster assignment}

The authors solve the shrinkage problem through an iterative reweighting procedure. At each iteration, the algorithm assigns weights to the residuals \(G-F\) and to differences between connected rows of \(G\). It then updates \(G\) by solving a linear system.

This procedure converts the nonsmooth objective into a sequence of more tractable weighted quadratic problems. The paper provides a convergence argument based on the decrease of the objective function during the iterations.

Once \(G\) has been obtained, K-means is applied to its rows to produce the final clustering assignments.

The distinction between these two stages is important. Convexity applies to the optimization of \(G\) when \(F\) and \(W\) are fixed. It does not make the preceding graph construction or the subsequent K-means optimization convex. Consequently, convergence of the shrinkage procedure does not, by itself, guarantee the best possible clustering accuracy.

SSC also connects naturally to several concepts covered in COMP3314. Like PCA and LDA-based clustering, it uses a lower-dimensional representation to make clustering more effective. Unlike PCA, which preserves directions of high variance, or discriminant approaches that exploit cluster structure, SSC begins with a graph-based spectral embedding. Its additional shrinkage step is an example of regularization: it changes the representation by balancing fidelity to existing information against a structural preference encoded by the neighborhood graph.

\section{Experimental Evidence in the Original Paper}

Chang et al.\ evaluate SSC on eight benchmarks: AR, JAFFE, ORL, UMIST, BinAlpha, MSRA50, YaleB, and USPS. These datasets include facial images, handwritten characters, and handwritten digits, providing several different clustering settings.

The comparison includes K-means, spectral clustering (SC), PCA followed by K-means or SC, LDA-based variants, Local Learning Clustering (LLC), and Subspace Learning via Pattern Shrinking (SLPS, referred to as SPLS in the SSC paper) [Ding and Li, 2007; Wu and Schölkopf, 2006; Hou et al., 2013].

The experiments use \(k=5\) for the neighborhood construction of spectral methods. The authors also describe a parameter search over \(\{10^{-6},10^{-3},1,10^3,10^6\}\), dimensionality-reduction experiments over 10--150 dimensions, and 50 randomized clustering repetitions. Results are evaluated using Clustering Accuracy (ACC) and Normalized Mutual Information (NMI).

The paper reports that SSC achieves the highest ACC across all eight datasets. For example, its ACC improves upon standard SC from 72.7\% to 76.0\% on ORL, 76.1\% to 81.2\% on JAFFE, 34.8\% to 49.8\% on YaleB, and 64.3\% to 75.5\% on USPS.

The authors also present a parameter-sensitivity study showing how clustering performance changes with \(\gamma\), together with convergence curves illustrating the decrease of the optimization objective.

These results support the paper's central empirical argument that modifying the spectral representation can improve clustering. However, not all reported metrics follow the same pattern. On YaleB, for instance, Table 2 reports an NMI of 46.8\% for SSC, compared with 55.6\% for standard SC, even though SSC achieves considerably higher ACC. Thus, the published evidence supports SSC's consistent ACC advantage more directly than an unconditional improvement in every evaluation metric.

\section{Independent Reproduction}

\subsection{Implementation and experimental protocol}

We reconstructed SSC from the mathematical formulation and algorithm described in the paper, without access to the original authors' implementation. Our implementation separates the initial spectral embedding, construction of the graph \(W\), shrinkage optimization, and final K-means assignment.

For the main comparison, we use the symmetric normalized Laplacian, \(k=5\), a union of directed nearest-neighbor relationships, and a Gaussian bandwidth determined by the median positive directed kNN squared distance. The spectral eigenvectors are used without row normalization.

These choices are explicitly documented because the paper does not fully specify the bandwidth calculation, graph symmetrization, or all details of the spectral discretization procedure.

The main SSC results use the objective-consistent implementation, denoted E8, with \(\gamma=1\), initialization \(G_0=F\), and a maximum of 300 optimization iterations. The final K-means stage uses random initialization with one initialization per trial and seeds 0--49. Each reported mean and standard deviation is calculated over 50 clustering trials on a fixed spectral representation.

This fixed-\(\gamma\) comparison is a controlled reference setting, not a reconstruction of the full parameter-selection procedure used for every published result. Alternative settings are investigated separately to avoid confusing diagnostic improvements with reproduction of the authors' historical protocol.

ACC is evaluated after optimal permutation of cluster labels using the Hungarian assignment algorithm [Munkres, 1957]. NMI follows the geometric-mean normalization of the two clustering entropies [Strehl and Ghosh, 2002].

\subsection{Benchmark availability}

The principal comparison uses four datasets for which suitable public source files were recovered:
\begin{itemize}
\item ORL: 400 AT\&T face images, resized from the original \(92\times112\) representation to \(32\times32\).
\item JAFFE: 213 facial-expression images from the archived JAFFE distribution, represented at \(26\times26\).
\item YaleB: 2,414 non-ambient Cropped YaleB images, resized to \(32\times32\).
\item USPS: 9,298 handwritten-digit samples obtained from the LIBSVM distribution.
\end{itemize}

These data are sufficient for a controlled partial reproduction, but matching the published dataset sizes does not establish that every historical preprocessing choice has been recovered.

Other benchmarks remain limited by source availability or sample identity. The recovered AR archive contains only 100 identities rather than the 120 required by the paper. Historical UMIST archives contain 575 cropped images, including duplicates, rather than an identifiable 564-image subset. The available Binary Alphadigits collection supports a 780-image reconstruction, but the original selection indices are unknown. The required MSRA50 data were not recovered.

Most baseline methods were independently implemented from their published descriptions. However, LDA-SC remains an approximate reconstruction because its exact procedure is not specified in sufficient detail, while SPLS could not be implemented faithfully without the full original algorithm. These methods are not presented as verified reproductions.

\subsection{Reproduction results}

Table 1 compares the published SC and SSC results with our independently measured results on the four available benchmark reconstructions. All values are percentages.

\input{tables/final_main}

The results reveal substantial differences across datasets. On ORL, our SSC implementation achieves 57.3\% ACC, compared with 56.2\% for SC. Although shrinkage provides a modest improvement, both results remain considerably below their published counterparts. The discrepancy therefore appears before the shrinkage optimization itself and cannot be explained entirely by the way \(G\) is computed.

A similar pattern appears on JAFFE, where SC and SSC produce nearly identical reproduced accuracy. On YaleB, the reproduced SSC accuracy is slightly below that of SC, despite the much larger improvement reported in the original paper. USPS shows almost no improvement at the fixed reference value \(\gamma=1\).

To investigate whether these differences arise from random K-means initialization, we compared clustering on the same \(F\) and \(G\) using matched sets of 50 seeds. On ORL, shrinkage increased mean ACC by approximately 1.09 percentage points, with a 95\% bootstrap confidence interval of 0.30--1.87 points. On YaleB, the corresponding change was approximately $-0.72$ points, with a confidence interval of $-0.97$ to $-0.48$ points. These intervals quantify variation due to K-means initialization on fixed representations; they do not account for uncertainty in dataset sampling or graph construction.

We also examined a second YaleB source, the commonly used Deng Cai \(32\times32\) feature matrix. Although this version and our resized Cropped YaleB data contain the same number of observations and classes, their pixel representations are not identical. Spectral clustering improves only modestly when the Deng Cai version is used, with ACC rising from approximately 18.3\% to 20.2\%. The remaining difference from the published 34.8\% suggests that changing the dataset version alone is insufficient to recover the original spectral-clustering result.

\section{Critical Analysis}

\subsection{Mathematical formulation and optimization}

A notable strength of SSC is that its representation-refinement objective has a clear convex structure. For fixed \(F\) and \(W\), the problem can be studied using standard tools from convex optimization, and small numerical instances can be checked against independent solvers.

However, the printed equations contain an ambiguity that matters for implementation. Under the convention in which both ordered pairs of a symmetric graph are included in the shrinkage sum, differentiating Equation (8) leads to a linear system of the form
\[
(S+2\gamma\widetilde L)G=SF,
\]
whereas Equation (13) uses \(S+\gamma\widetilde L\). Algorithm 1 also contains inconsistent notation in its update step.

We therefore distinguish the objective-consistent reconstruction, E8, from an implementation following the printed Equation (13), E13. On a six-point numerical test, the E8 result agrees with an independent solution of the printed objective to within approximately \(1.17\times10^{-5}\) in the maximum absolute coordinate difference. The printed-update interpretation agrees instead with a formulation using a corresponding rescaling of \(\gamma\).

This difference is important for scientific reproducibility, although it does not establish which convention was used in the original authors' code. Nor does agreement on a small optimization instance guarantee that every large-scale run has converged sufficiently.

More fundamentally, the convexity result concerns the learned representation, not the entire clustering procedure. The affinity graphs depend on the data, and the final K-means objective remains nonconvex. SSC may therefore solve its representation-learning subproblem accurately without producing an equally accurate discrete partition.

\subsection{Sensitivity of the clustering pipeline}

Our experiments also highlight the dependence of SSC on the initial spectral embedding. The shrinkage procedure is applied after \(F\) has already been constructed, so errors in the initial graph or spectral representation may persist into the subsequent optimization.

This issue is particularly visible on ORL and YaleB, where standard SC already performs substantially below the published values. The additional shrinkage step cannot be treated as an independent remedy for all errors introduced earlier in the pipeline.

A further source of variation is the final K-means initialization. The original paper reports repeated random initializations and objective-based selection, but the exact relationship between this selection rule and the published mean $\pm$ standard deviation is not fully specified. Our controlled results show that differences in clustering assignments can arise even when the learned representation is held fixed.

Finally, the amount of shrinkage is not synonymous with clustering quality. Increasing \(\gamma\) can produce a smoother representation, but greater geometric contraction does not necessarily preserve or strengthen the boundaries between meaningful clusters. Our results show that the accuracy gained from shrinkage is small or negative in several of the reconstructed benchmark settings.

\subsection{Reproducibility limitations}

The remaining discrepancies should be interpreted in light of several limitations. The original code is unavailable, some historical benchmark subsets cannot be identified precisely, and implementation choices such as Gaussian bandwidth selection and graph symmetrization are not completely specified.

The reported parameter-selection procedure also leaves room for different interpretations. These uncertainties matter because the method combines several individually sensitive stages: graph construction, spectral embedding, shrinkage optimization, and K-means discretization.

Our experiments therefore do not disprove the results reported by Chang et al. Rather, they establish that the published improvements have not been consistently recovered under our documented, independently implemented protocol.

This distinction is central to the assessment of reproducibility. A method can be mathematically well defined in its main objective while still requiring substantial undocumented experimental detail to reproduce its published numerical performance.

\section{Insights and proposed extension}

\section{Conclusion}

Spectral Shrunk Clustering offers a principled way to refine spectral embeddings by encouraging neighboring points to move together. Its convex formulation is an attractive feature, and the method provides an instructive connection between graph-based representation learning, regularization, and clustering.

Our independent implementation supports the mathematical interpretation of the shrinkage objective but does not consistently reproduce the published gains over standard spectral clustering. The results indicate that the success of SSC depends not only on its convex optimization step, but also on the quality of the initial embedding, the construction of its neighborhood graph, and the final clustering procedure.

The reproduction remains partial because of unavailable historical resources and unresolved implementation details. Nevertheless, the analysis demonstrates why mathematical correctness and experimental reproducibility must be evaluated separately, particularly in multi-stage unsupervised learning methods.

\begin{thebibliography}{18}

\bibitem{chang2015}
Chang, X., Nie, F., Ma, Z., Yang, Y., \& Zhou, X. (2015).
A convex formulation for spectral shrunk clustering.
Proceedings of the Twenty-Ninth AAAI Conference on Artificial Intelligence, 2532--2538.
\url{https://doi.org/10.1609/aaai.v29i1.9606}

\bibitem{ding2007}
Ding, C., \& Li, T. (2007).
Adaptive dimension reduction using discriminant analysis and K-means clustering.
Proceedings of the 24th International Conference on Machine Learning, 521--528.
\url{https://doi.org/10.1145/1273496.1273562}

\bibitem{georghiades2001}
Georghiades, A. S., Belhumeur, P. N., \& Kriegman, D. J. (2001).
From few to many: Illumination cone models for face recognition under variable lighting and pose.
IEEE Transactions on Pattern Analysis and Machine Intelligence, 23(6), 643--660.
\url{https://doi.org/10.1109/34.927464}

\bibitem{graham1998}
Graham, D. B., \& Allinson, N. M. (1998).
Characterising virtual eigensignatures for general purpose face recognition.
In Face Recognition: From Theory to Applications, 446--456. Springer.
\url{https://doi.org/10.1007/978-3-642-72201-1_25}

\bibitem{hou2013}
Hou, C., Nie, F., Jiao, Y., Zhang, C., \& Wu, Y. (2013).
Learning a subspace for clustering via pattern shrinking.
Information Processing \& Management, 49(4), 871--883.
\url{https://doi.org/10.1016/j.ipm.2013.01.004}

\bibitem{hull1994}
Hull, J. J. (1994).
A database for handwritten text recognition research.
IEEE Transactions on Pattern Analysis and Machine Intelligence, 16(5), 550--554.
\url{https://doi.org/10.1109/34.291440}

\bibitem{lee2005}
Lee, K.-C., Ho, J., \& Kriegman, D. J. (2005).
Acquiring linear subspaces for face recognition under variable lighting.
IEEE Transactions on Pattern Analysis and Machine Intelligence, 27(5), 684--698.
\url{https://doi.org/10.1109/TPAMI.2005.92}

\bibitem{lloyd1982}
Lloyd, S. P. (1982).
Least squares quantization in PCM.
IEEE Transactions on Information Theory, 28(2), 129--137.
\url{https://doi.org/10.1109/TIT.1982.1056489}

\bibitem{lyons1998}
Lyons, M., Akamatsu, S., Kamachi, M., \& Gyoba, J. (1998).
Coding facial expressions with Gabor wavelets.
Proceedings of the Third IEEE International Conference on Automatic Face and Gesture Recognition, 200--205.
\url{https://doi.org/10.1109/AFGR.1998.670949}

\bibitem{lyons2025}
Lyons, M. J., Kamachi, M., \& Gyoba, J. (2025).
The Japanese Female Facial Expression (JAFFE) Dataset.
Zenodo, Record 14974867.

\bibitem{munkres1957}
Munkres, J. (1957).
Algorithms for the assignment and transportation problems.
Journal of the Society for Industrial and Applied Mathematics, 5(1), 32--38.
\url{https://doi.org/10.1137/0105003}

\bibitem{ng2001}
Ng, A. Y., Jordan, M. I., \& Weiss, Y. (2001).
On spectral clustering: Analysis and an algorithm.
Advances in Neural Information Processing Systems, 14, 849--856.

\bibitem{nie2010}
Nie, F., Huang, H., Cai, X., \& Ding, C. (2010).
Efficient and robust feature selection via joint $\ell_{2,1}$-norms minimization.
Advances in Neural Information Processing Systems, 23, 1813--1821.

\bibitem{samaria1994}
Samaria, F. S., \& Harter, A. C. (1994).
Parameterisation of a stochastic model for human face identification.
Proceedings of the IEEE Workshop on Applications of Computer Vision, 138--142.
\url{https://doi.org/10.1109/ACV.1994.341300}

\bibitem{shi2000}
Shi, J., \& Malik, J. (2000).
Normalized cuts and image segmentation.
IEEE Transactions on Pattern Analysis and Machine Intelligence, 22(8), 888--905.
\url{https://doi.org/10.1109/34.868688}

\bibitem{strehl2002}
Strehl, A., \& Ghosh, J. (2002).
Cluster ensembles---A knowledge reuse framework for combining multiple partitions.
Journal of Machine Learning Research, 3, 583--617.

\bibitem{vonluxburg2007}
von Luxburg, U. (2007).
A tutorial on spectral clustering.
Statistics and Computing, 17(4), 395--416.
\url{https://doi.org/10.1007/s11222-007-9033-z}

\bibitem{wu2006}
Wu, M., \& Schölkopf, B. (2006).
A local learning approach for clustering.
Advances in Neural Information Processing Systems, 19, 1529--1536.

\end{thebibliography}
\makeatletter
\global\NAT@stdbstfalse
\makeatother

\end{document}
